
\documentclass[11pt,a4paper]{article}
\usepackage[T1]{fontenc}
\usepackage{lmodern}
\usepackage{amsmath,amssymb,amsfonts,mathtools}
\usepackage{graphicx,booktabs,longtable,array,tabularx,adjustbox}
\usepackage{geometry,microtype,fancyhdr,lastpage,etoolbox,enumitem}
\usepackage{newunicodechar,textcomp}
\usepackage{hyperref}
\hypersetup{hidelinks,pdfcreator={XeLaTeX / QM1 academic edition}}
\geometry{left=20mm,right=20mm,top=22mm,bottom=23mm}
\setcounter{secnumdepth}{2}
\setlength{\parindent}{0pt}
\setlength{\parskip}{4pt plus 1pt}
\setlength{\emergencystretch}{2em}
\setlength{\tabcolsep}{3pt}
\renewcommand{\arraystretch}{1.16}
\setlist{nosep,topsep=3pt,itemsep=2pt,leftmargin=7mm}
\AtBeginEnvironment{longtable}{\footnotesize\setlength{\tabcolsep}{2.2pt}}
\newcommand{\tightlist}{\setlength{\itemsep}{0pt}\setlength{\parskip}{0pt}}
\providecommand{\Cov}{\operatorname{Cov}}
\providecommand{\Var}{\operatorname{Var}}
\providecommand{\Corr}{\operatorname{Corr}}
\providecommand{\MAD}{\operatorname{MAD}}
\providecommand{\E}{\mathbb{E}}
\providecommand{\R}{\mathbb{R}}
\newunicodechar{—}{\textemdash}
\newunicodechar{–}{\textendash}
\newunicodechar{€}{\texteuro}
\newunicodechar{≤}{\ensuremath{\leq}}
\newunicodechar{≥}{\ensuremath{\geq}}
\newunicodechar{×}{\ensuremath{\times}}
\newunicodechar{÷}{\ensuremath{\div}}
\newunicodechar{Σ}{\ensuremath{\Sigma}}
\newunicodechar{σ}{\ensuremath{\sigma}}
\newunicodechar{μ}{\ensuremath{\mu}}
\newunicodechar{Δ}{\ensuremath{\Delta}}
\newunicodechar{∑}{\ensuremath{\sum}}
\newunicodechar{≈}{\ensuremath{\approx}}
\newunicodechar{√}{\ensuremath{\surd}}
\newunicodechar{≠}{\ensuremath{\ne}}
\newunicodechar{→}{\ensuremath{\to}}
\newunicodechar{∞}{\ensuremath{\infty}}
\newunicodechar{∈}{\ensuremath{\in}}
\newunicodechar{−}{\ensuremath{-}}
\newunicodechar{²}{\ensuremath{^{2}}}
\newunicodechar{³}{\ensuremath{^{3}}}
\newunicodechar{₀}{\ensuremath{_{0}}}
\newunicodechar{₁}{\ensuremath{_{1}}}
\newunicodechar{₂}{\ensuremath{_{2}}}
\newunicodechar{₃}{\ensuremath{_{3}}}
\newunicodechar{₄}{\ensuremath{_{4}}}
\newunicodechar{ₜ}{\ensuremath{_{t}}}
\newunicodechar{ₚ}{\ensuremath{_{p}}}
\newunicodechar{ₙ}{\ensuremath{_{n}}}
\newunicodechar{ₓ}{\ensuremath{_{x}}}
\newunicodechar{ᵢ}{\ensuremath{_{i}}}
\newunicodechar{ŷ}{\ensuremath{\hat y}}
\newunicodechar{Ŷ}{\ensuremath{\widehat Y}}
\newunicodechar{ₘ}{\ensuremath{_{m}}}
\newunicodechar{ₑ}{\ensuremath{_{e}}}
\newunicodechar{ₐ}{\ensuremath{_{a}}}
\newunicodechar{ₖ}{\ensuremath{_{k}}}
\newunicodechar{ₛ}{\ensuremath{_{s}}}
\newunicodechar{ₗ}{\ensuremath{_{l}}}
\newunicodechar{ₕ}{\ensuremath{_{h}}}
\newunicodechar{ₒ}{\ensuremath{_{o}}}
\newunicodechar{₊}{\ensuremath{_{+}}}
\newunicodechar{₋}{\ensuremath{_{-}}}
\newunicodechar{₌}{\ensuremath{_{=}}}
\newunicodechar{₍}{\ensuremath{_{(}}}
\newunicodechar{₎}{\ensuremath{_{)}}}
\newunicodechar{⁻}{\ensuremath{^{-}}}
\newunicodechar{⁺}{\ensuremath{^{+}}}
\newunicodechar{⁰}{\ensuremath{^{0}}}
\newunicodechar{¹}{\ensuremath{^{1}}}
\newunicodechar{⁴}{\ensuremath{^{4}}}
\newunicodechar{⁵}{\ensuremath{^{5}}}
\newunicodechar{⁶}{\ensuremath{^{6}}}
\newunicodechar{⁷}{\ensuremath{^{7}}}
\newunicodechar{⁸}{\ensuremath{^{8}}}
\newunicodechar{⁹}{\ensuremath{^{9}}}
\makeatletter
\renewcommand{\section}{\@startsection{section}{1}{0pt}{2.0ex plus .7ex minus .3ex}{1.0ex}{\large\bfseries}}
\renewcommand{\subsection}{\@startsection{subsection}{2}{0pt}{1.8ex plus .5ex minus .2ex}{.7ex}{\normalsize\bfseries}}
\renewcommand{\subsubsection}{\@startsection{subsubsection}{3}{0pt}{1.3ex plus .4ex minus .2ex}{.5ex}{\normalsize\bfseries}}
\makeatother
\pagestyle{fancy}
\fancyhf{}
\renewcommand{\headrulewidth}{0.35pt}
\renewcommand{\footrulewidth}{0pt}
\fancyhead[L]{\small\textsc{KEDGE Business School --- QM1}}
\fancyhead[R]{\small\textsc{2026--2027}}
\fancyfoot[L]{\small Badr ABOUFARASSE, PhD}
\fancyfoot[R]{\small\thepage\ /\ \pageref{LastPage}}
\setlength{\headheight}{14pt}
\setkeys{Gin}{keepaspectratio}
\begin{document}

\begin{center}
{\small\textsc{KEDGE Business School \quad EBP1 \quad Quantitative Methods 1}}\\[6pt]
{\Large\bfseries Worksheet 3 — Dispersion and concentration}\\[6pt]
{\normalsize\bfseries Exercise statements}\\[6pt]
{\small Badr ABOUFARASSE, PhD \quad --- \quad 2026--2027}
\end{center}
\vspace{2mm}\hrule\vspace{6mm}

\hypertarget{core-applications}{%
\section{Core applications}\label{core-applications}}

Master the essential definitions and calculations from the course.

\emph{Core application --- 25 min}

\hypertarget{exercise-4-medial-and-cumulative-relative-masses}{%
\subsection{Exercise 4 --- Medial and cumulative relative masses}\label{exercise-4-medial-and-cumulative-relative-masses}}

The spending of 40 customers (€) is grouped below. Approximate class monetary mass by nᵢ × class midpoint.

\begin{longtable}[]{@{}llllll@{}}
\toprule\noalign{}
Classe / Class & {[}10,20{[} & {[}20,30{[} & {[}30,40{[} & {[}40,50{[} & {[}50,60{[} \\
\midrule\noalign{}
\endhead
\bottomrule\noalign{}
\endlastfoot
nᵢ & 8 & 12 & 10 & 6 & 4 \\
\end{longtable}

\begin{enumerate}
\tightlist
\item
  Compute midpoints, masses, cumulative masses and relative cumulative mass shares.
\item
  Interpolate the medial in the class containing half of total mass.
\item
  Compute the frequency median and compare the two statistics.
\end{enumerate}

\emph{Core application --- 20 min}

\hypertarget{extra-1-dispersion-from-a-frequency-table}{%
\subsection{Extra 1 --- Dispersion from a frequency table}\label{extra-1-dispersion-from-a-frequency-table}}

The number of weekly complaints per store is distributed as follows:

\begin{longtable}[]{@{}llllll@{}}
\toprule\noalign{}
Complaints x & 0 & 1 & 2 & 3 & 4 \\
\midrule\noalign{}
\endhead
\bottomrule\noalign{}
\endlastfoot
Stores n & 4 & 8 & 10 & 6 & 2 \\
\end{longtable}

\begin{enumerate}
\tightlist
\item
  Calculate the mean, population variance and standard deviation.
\item
  Calculate the coefficient of variation.
\item
  Interpret every result in context.
\end{enumerate}

\textbf{Method hint:}

Build columns nx and nx². Use E(X²) − {[}E(X){]}² for the population variance.

\hypertarget{assessment-and-exam-practice}{%
\section{Assessment and exam practice}\label{assessment-and-exam-practice}}

Practise independently with assessment and final-exam style questions.

\emph{Assessment / exam practice --- 35 min}

\hypertarget{exercise-2-dispersion-of-grouped-data}{%
\subsection{Exercise 2 -- Dispersion of grouped data}\label{exercise-2-dispersion-of-grouped-data}}

Using the 210-customer distribution from Worksheet 2, calculate {\(\sum n_ic_i\)}, {\(\sum n_ic_i^2\)}, mean, population variance, standard deviation, coefficient of variation and IQR. Compare the information provided by standard deviation and IQR.

\emph{Assessment / exam practice --- 30 min}

\hypertarget{exercise-3-lorenz-curve-and-gini-coefficient}{%
\subsection{Exercise 3 -- Lorenz curve and Gini coefficient}\label{exercise-3-lorenz-curve-and-gini-coefficient}}

Income values are {\(10,12,18,25,35\)}. Calculate cumulative population and income shares, construct the Lorenz curve, calculate the area by trapezoids, derive the Gini coefficient and explain why Gini measures concentration without explaining its causes.

\hypertarget{optional-advanced-practice}{%
\section{Optional advanced practice}\label{optional-advanced-practice}}

Explore additional interpretations and professional applications.

\emph{Advanced practice --- 30 min}

\hypertarget{exercise-1-boxplot-potential-outliers-and-variance-choice}{%
\subsection{Exercise 1 -- Boxplot, potential outliers and variance choice}\label{exercise-1-boxplot-potential-outliers-and-variance-choice}}

Consider {\[5,6,6,7,7,8,8,9,9,10,11,20.\]} Use the median-of-halves convention for quartiles.

\begin{enumerate}
\item
  Calculate the five-number summary.
\item
  Calculate the range and IQR.
\item
  Calculate the fences {\(Q_1-1.5IQR\)} and {\(Q_3+1.5IQR\)} and identify potential outliers.
\item
  Calculate the mean and mean absolute deviation.
\item
  Calculate population variance and standard deviation if the 12 observations are the entire population.
\item
  Calculate sample variance and standard deviation if they are a sample.
\end{enumerate}

\emph{Advanced practice --- 25 min}

\hypertarget{extra-2-comparing-concentration-and-robustness}{%
\subsection{Extra 2 --- Comparing concentration and robustness}\label{extra-2-comparing-concentration-and-robustness}}

Two markets have the following ordered firm market shares:

\begin{longtable}[]{@{}llllll@{}}
\toprule\noalign{}
Market A & 10\% & 15\% & 20\% & 25\% & 30\% \\
Market B & 2\% & 8\% & 10\% & 20\% & 60\% \\
\midrule\noalign{}
\endhead
\bottomrule\noalign{}
\endlastfoot
\end{longtable}

\begin{enumerate}
\tightlist
\item
  Construct cumulative shares for each market.
\item
  Explain which Lorenz curve is farther from equality.
\item
  For each market calculate the range and discuss why it is insufficient.
\item
  State a defensible concentration conclusion.
\end{enumerate}

\textbf{Method hint:}

With five equal-weight firms, cumulative population shares increase by 20 percentage points.

\end{document}
